\chapter{Ерөнхий дүгнэлт}

\section{Хийгдсэн ажлын нэгтгэл}
Энэхүү төслийн хүрээнд нээлттэй эхийн Fossify Messages
аппликейшнд суурилан дараах үндсэн бүрэлдэхүүн хэсгүүдийг
амжилттай хөгжүүлж, хэрэгжүүллээ:

\begin{itemize}
    \item \textbf{Тогтмол илэрхийллийн шүүлтүүр}:
    Хэрэглэгчийн тодорхойлсон Regex болон түлхүүр үгэнд
    суурилсан автомат ангиллын алгоритм хэрэгжүүлсэн.

    \item \textbf{Хадгалсан харагдацын систем}:
    Хэрэглэгч өөрийн хэрэгцээнд тааруулсан хавтас
    үүсгэж, зурвасаа ангилан харах боломжийг бүрдүүлсэн.

    \item \textbf{Нугалдаг төхөөрөмжийн дэмжлэг}:
    Samsung Galaxy Z Fold3 дээр хоёр баганат харагдацыг
    гацалтгүйгээр ажиллуулсан.

    \item \textbf{Шудрах үйлдлийн тохиргоо}:
    Хэрэглэгч архивлах, устгах, уншсан тэмдэглэх зэрэг
    үйлдлийг шудрах хөдөлгөөнд оноох боломжтой болсон.

    \item \textbf{Хугацааны бүлэглэлт}:
    Inbox дахь харилцан яриануудыг өнөөдөр, өчигдөр,
    энэ сар гэсэн бүлгүүдэд автоматаар ангилдаг болсон.
\end{itemize}

\section{Зорилгын биелэлт}
Chapter 1-д SMART зарчмаар тавьсан зорилгуудын биелэлтийг
дараах байдлаар нэгтгэв:

\begin{table}[H]
    \centering
    \caption{SMART зорилгын биелэлтийн нэгтгэл}
    \label{table:smart_result}
    \begin{tabular}{|l|p{5cm}|c|}
        \hline
        \textbf{Зорилго} & \textbf{Үр дүн} & \textbf{Биелэлт} \\
        \hline
        Specific & Тогтмол илэрхийлэл болон түлхүүр үгэнд суурилсан
        автомат ангилал хэрэгжсэн & \checkmark \\
        \hline
        Measurable & 31 тест 100\% Pass,
        гол логикийн coverage ~85\% & \checkmark \\
        \hline
        Achievable & Нугалдаг төхөөрөмж дээр
        тасралтгүй 2 баганат харагдац & \checkmark \\
        \hline
        Relevant & Нээлттэй эхийн орчинд
        хэрэгжүүлсэн & \checkmark \\
        \hline
        Time-bound & Хамгаалалтын өмнө
        бүрэн ажиллагаатай & \checkmark \\
        \hline
    \end{tabular}
\end{table}

\section{Туршилтын үр дүнгийн нэгтгэл}
Системийн найдвартай ажиллагааг баталгаажуулахын тулд
нэгж туршилт, нэгтгэл туршилт болон ачааллын туршилт
гүйцэтгэсэн үр дүнг дараах байдлаар нэгтгэв:

\begin{itemize}
    \item \textbf{Нэгж туршилт}: 31 тест, 100\% Pass
    \item \textbf{Цөм логикийн кодын хамралт}: ~85\%
    \item \textbf{Ачааллын туршилт}: 40,000 харилцан яриа,
    дунджаар 15мс-ээс бага хариу хугацаа
    \item \textbf{Бодит төхөөрөмж}: Samsung Galaxy Note20 Ultra
    болон Z Fold3 дээр бүх үндсэн функцүүд амжилттай ажилласан
\end{itemize}

\section{Хязгаарлалт}
Төслийн хүрээнд дараах зүйлсийг хэрэгжүүлэх боломжгүй
байсан тул цаашдын ажил болгон үлдээв:

\begin{itemize}
    \item CI/CD pipeline-г Fossify-ийн нарийн
    build dependency-аас шалтгаалан бүрэн автоматжуулж
    чадаагүй
    \item Fossify Messages repository-д өөрийн хөгжүүлсэн функцуудыг нэгтгээгүй.
\end{itemize}

\section{Цаашдын хөгжүүлэлт}
Энэхүү төсөл нь цаашид илүү өргөжүүлэн хөгжүүлж,
хэрэглэгчдийн өдөр тутмын амьдралыг илүү хялбаршуулах
бүрэн боломжтой юм. Тэргүүлэх чиглэлүүд:

\begin{itemize}
    \item \textbf{Шүүх механизмын сайжруулалт}:
    Ангилалах дүрэмд нэмэлт нөхцөлүүдийг (жишээ нь: Харилцагчын бүлэг, ангилагдахгүй нөхцөл, цаг хугацаа гэх мэт) нэмж, илүү нарийн ангилал хийх боломжийг бий болгох
    \item \textbf{CI/CD бүрэн нэвтрүүлэлт}:
    GitHub Actions pipeline-г тогтворжуулж,
    автомат тест болон APK build-ийг хэрэгжүүлэх
    \item \textbf{Агуулгад суурилсан функцууд}:
    Ангиллын дагууд мэдэгдэл тохируулах боломжийг нэмэх, агуулгыг танаж багасгах
    \item \textbf{Fossify Messages-т нэгтгэх}:
    Нийтэд нээлттэй репозиторт өөрийн хөгжүүлсэн кодыг нэгтгэж,
    нээлттэй эх бүхий программ хангамж хөгжүүлэгчдийн нийгэмлэгийн хамтын ажиллагаанд нэмэр болох
\end{itemize}

Өдөр тутмын амьдралд бидэнд бэлэн өгөгддөг зүйлсийг
өөрийн хэрэгцээ шаардлагын дагуу өөрчилж, сайжруулснаар
ямар их эерэг өөрчлөлтийг авчирч болдгийг энэхүү ажил харууллаа.
